\documentclass[10pt,a4paper]{amsart}
\usepackage{amsmath,amssymb,amsthm,graphicx}
\usepackage{verbatim}
\newenvironment{reference}{\comment}{\endcomment}
%\newenvironment{reference}{}{}
\newenvironment{slogan}{\comment}{\endcomment}
\newenvironment{history}{\comment}{\endcomment}
% For commutative diagrams we use Xy-pic
\usepackage[all]{xy}

% We use 2cell for 2-commutative diagrams.
\xyoption{2cell}
% Theorem environments.
%
\theoremstyle{plain}
\newtheorem{theorem}[subsection]{Theorem}
\newtheorem{proposition}[subsection]{Proposition}
\newtheorem{lemma}[subsection]{Lemma}

\theoremstyle{definition}
\newtheorem{definition}[subsection]{Definition}
\newtheorem{example}[subsection]{Example}
\newtheorem{exercise}[subsection]{Exercise}
\newtheorem{situation}[subsection]{Situation}

\theoremstyle{remark}
\newtheorem{remark}[subsection]{Remark}
\newtheorem{remarks}[subsection]{Remarks}

\numberwithin{equation}{subsection}

% Macros
%
\def\lim{\mathop{\mathrm{lim}}\nolimits}
\def\colim{\mathop{\mathrm{colim}}\nolimits}
\def\Spec{\mathop{\mathrm{Spec}}}
\def\Hom{\mathop{\mathrm{Hom}}\nolimits}
\def\Ext{\mathop{\mathrm{Ext}}\nolimits}
\def\SheafHom{\mathop{\mathcal{H}\!\mathit{om}}\nolimits}
\def\SheafExt{\mathop{\mathcal{E}\!\mathit{xt}}\nolimits}
\def\Sch{\mathit{Sch}}
\def\Mor{\mathop{\mathrm{Mor}}\nolimits}
\def\Ob{\mathop{\mathrm{Ob}}\nolimits}
\def\Sh{\mathop{\mathit{Sh}}\nolimits}
\def\NL{\mathop{N\!L}\nolimits}
\def\CH{\mathop{\mathrm{CH}}\nolimits}
\def\proetale{{pro\text{-}\acute{e}tale}}
\def\etale{{\acute{e}tale}}
\def\QCoh{\mathit{QCoh}}
\def\Ker{\mathop{\mathrm{Ker}}}
\def\Im{\mathop{\mathrm{Im}}}
\def\Coker{\mathop{\mathrm{Coker}}}
\def\Coim{\mathop{\mathrm{Coim}}}

% Boxtimes
%
\DeclareMathSymbol{\boxtimes}{\mathbin}{AMSa}{"02}

%
% Macros for moduli stacks/spaces
%
\def\QCohstack{\mathcal{QC}\!\mathit{oh}}
\def\Cohstack{\mathcal{C}\!\mathit{oh}}
\def\Spacesstack{\mathcal{S}\!\mathit{paces}}
\def\Quotfunctor{\mathrm{Quot}}
\def\Hilbfunctor{\mathrm{Hilb}}
\def\Curvesstack{\mathcal{C}\!\mathit{urves}}
\def\Polarizedstack{\mathcal{P}\!\mathit{olarized}}
\def\Complexesstack{\mathcal{C}\!\mathit{omplexes}}
% \Pic is the operator that assigns to X its picard group, usage \Pic(X)
% \Picardstack_{X/B} denotes the Picard stack of X over B
% \Picardfunctor_{X/B} denotes the Picard functor of X over B
\def\Pic{\mathop{\mathrm{Pic}}\nolimits}
\def\Picardstack{\mathcal{P}\!\mathit{ic}}
\def\Picardfunctor{\mathrm{Pic}}
\def\Deformationcategory{\mathcal{D}\!\mathit{ef}}
\usepackage{cleveref}
\crefname{originallemma}{original lemma}{original lemmas}
\Crefname{originallemma}{Original lemma}{Original lemmas}
\makeatletter
\newlabel{lemma-bivariant-zero}{{02UC}{source}{}{sourcepointer-1}{}}
\newlabel{lemma-bivariant-zero@cref}{{[originallemma][0][]02UC}{[1][0][]source}{}{}{}}
\makeatother
\crefname{originallemma}{original lemma}{original lemmas}
\Crefname{originallemma}{Original lemma}{Original lemmas}
\makeatletter
\newlabel{lemma-determine-intersections}{{02U6}{source}{}{sourcepointer-2}{}}
\newlabel{lemma-determine-intersections@cref}{{[originallemma][0][]02U6}{[1][0][]source}{}{}{}}
\makeatother
\crefname{originallemma}{original lemma}{original lemmas}
\Crefname{originallemma}{Original lemma}{Original lemmas}
\makeatletter
\newlabel{lemma-pushforward-cap-c1}{{02SU}{source}{}{sourcepointer-3}{}}
\newlabel{lemma-pushforward-cap-c1@cref}{{[originallemma][0][]02SU}{[1][0][]source}{}{}{}}
\makeatother
\crefname{originallemma}{original lemma}{original lemmas}
\Crefname{originallemma}{Original lemma}{Original lemmas}
\makeatletter
\newlabel{lemma-relative-effective-cartier}{{02TR}{source}{}{sourcepointer-4}{}}
\newlabel{lemma-relative-effective-cartier@cref}{{[originallemma][0][]02TR}{[1][0][]source}{}{}{}}
\makeatother
\crefname{originalreference}{original reference}{original references}
\Crefname{originalreference}{Original reference}{Original references}
\makeatletter
\newlabel{situation-setup}{{02QL}{source}{}{sourcepointer-5}{}}
\newlabel{situation-setup@cref}{{[originalreference][0][]02QL}{[1][0][]source}{}{}{}}
\makeatother
\title{Chern classes of a finite locally free sheaf with a trivial line subbundle as a quotient}
\author{The Stacks Project Authors}
\date{stacks-project:chow.tex:local-snapshot-20261012; GNU FDL 1.2 or later}
\begin{document}
\maketitle
\noindent\textit{Editorial note.} All delivered mathematical text is the author's own, copied from the pinned Stacks Project chapter file chow.tex; no statement or proof step was rewritten, reordered or removed.\par\medskip
\noindent\textit{Editorial note.} Preamble provenance: the pinned chapter file opens with \textbackslash{}input\{preamble\} and carries no preamble text of its own, so the project's single Stacks-wide preamble staged at sources/stacks-chow/preamble.tex is excerpted verbatim as three spans, lines 8-12 (the verbatim package together with the project's own reference, slogan and history display-environment definitions, including the author's commented-out older form of the reference environment), lines 14-18 (Xy-pic setup) and lines 154-217 (the native theorem declarations and mathematical macros). The lines 8-12 span is required because a counted statement or its surrounding text may use those display environments, which the project typesets as comments. Nothing was added, removed or redefined; the wrapper class and the amsmath/amssymb/amsthm/graphicx base packages are supplied by the assembler.\par\medskip
\noindent\textit{Editorial note.} Same-chapter pointers are located at their \textbackslash{}label lines inside the pinned chapter file; cross-chapter pointers are located in the project's own tag list staged in the snapshot (sources/stacks-chow/tags.txt), at the verbatim lines carrying their canonical tags.\par\medskip
\noindent\textit{Allowed standard prerequisites (editorial).}\begin{itemize}
\item Vanishing criterion for bivariant classes. Same chapter chow.tex, label lemma-bivariant-zero, invoked at proof line 7881. Reduces the equality of bivariant classes to the equality after capping with the fundamental class of an integral base, which is the first reduction of the proof.
\item The cycle class of an effective Cartier divisor with trivial normal bundle. Same chapter chow.tex, label lemma-relative-effective-cartier, invoked at proof line 7902. Identifies the cap of the first Chern class of the twisting sheaf with the pushforward of the pullback along the closed immersion P' -> P, because P' is cut out by a section and O\_P(P') is the twisting sheaf.
\item Characterisation of the Chern classes by one relation on the projective bundle, and the projection formula. Same chapter chow.tex, labels lemma-determine-intersections (proof lines 7903 and 7917) and lemma-pushforward-cap-c1 (proof line 7911). The first gives both the relation satisfied by the classes of F and the uniqueness that identifies the classes of E from the pushed forward relation, and the second moves the cap with a power of the first Chern class through the pushforward.
\end{itemize}\medskip
\ifdefined\hypertarget\hypertarget{sourcepointer-1}{}\fi
\noindent\textit{Source reference.} Original lemma 02UC: sources/stacks-chow/pinned-original-source.tex, source lines 6585--6585. Navigation pointer for \textbackslash{}ref\{lemma-bivariant-zero\} used inside the counted span (line 7881); the argument that consumes it is retained in the span.\par\medskip
\ifdefined\hypertarget\hypertarget{sourcepointer-2}{}\fi
\noindent\textit{Source reference.} Original lemma 02U6: sources/stacks-chow/pinned-original-source.tex, source lines 7280--7280. Navigation pointer for \textbackslash{}ref\{lemma-determine-intersections\} used inside the counted span (line 7903); the argument that consumes it is retained in the span.\par\medskip
\ifdefined\hypertarget\hypertarget{sourcepointer-3}{}\fi
\noindent\textit{Source reference.} Original lemma 02SU: sources/stacks-chow/pinned-original-source.tex, source lines 4631--4631. Navigation pointer for \textbackslash{}ref\{lemma-pushforward-cap-c1\} used inside the counted span (line 7911); the argument that consumes it is retained in the span.\par\medskip
\ifdefined\hypertarget\hypertarget{sourcepointer-4}{}\fi
\noindent\textit{Source reference.} Original lemma 02TR: sources/stacks-chow/pinned-original-source.tex, source lines 5598--5598. Navigation pointer for \textbackslash{}ref\{lemma-relative-effective-cartier\} used inside the counted span (line 7902); the argument that consumes it is retained in the span.\par\medskip
\ifdefined\hypertarget\hypertarget{sourcepointer-5}{}\fi
\noindent\textit{Source reference.} Original reference 02QL: sources/stacks-chow/pinned-original-source.tex, source lines 1372--1372. Navigation pointer for \textbackslash{}ref\{situation-setup\} used inside the counted span (line 7864); the argument that consumes it is retained in the span.\par\medskip
\begin{lemma}
\label{lemma-get-rid-of-trivial-subbundle}
Let $(S, \delta)$ be as in Situation \ref{situation-setup}.
Let $X$ be locally of finite type over $S$.
Let $\mathcal{E}$, $\mathcal{F}$ be finite locally free sheaves
on $X$ of ranks $r$, $r - 1$ which fit into a short
exact sequence
$$
0 \to \mathcal{O}_X \to \mathcal{E} \to \mathcal{F} \to 0
$$
Then we have
$$
c_r(\mathcal{E}) = 0, \quad
c_j(\mathcal{E}) = c_j(\mathcal{F}), \quad j = 0, \ldots, r - 1
$$
in $A^*(X)$.
\end{lemma}

\begin{proof}
By Lemma \ref{lemma-bivariant-zero}
it suffices to show that if $X$ is integral
then $c_j(\mathcal{E}) \cap [X] = c_j(\mathcal{F}) \cap [X]$.
Let $(\pi : P \to X, \mathcal{O}_P(1))$,
resp.\ $(\pi' : P' \to X, \mathcal{O}_{P'}(1))$ denote the
projective space bundle associated to $\mathcal{E}$, resp.\ $\mathcal{F}$.
The surjection $\mathcal{E} \to \mathcal{F}$ gives rise
to a closed immersion
$$
i : P' \longrightarrow P
$$
over $X$. Moreover, the element
$1 \in \Gamma(X, \mathcal{O}_X) \subset \Gamma(X, \mathcal{E})$
gives rise to a global section $s \in \Gamma(P, \mathcal{O}_P(1))$
whose zero set is exactly $P'$. Hence $P'$ is an effective Cartier
divisor on $P$ such that $\mathcal{O}_P(P') \cong \mathcal{O}_P(1)$.
Hence we see that
$$
c_1(\mathcal{O}_P(1)) \cap \pi^*\alpha = i_*((\pi')^*\alpha)
$$
for any cycle class $\alpha$ on $X$ by
Lemma \ref{lemma-relative-effective-cartier}.
By Lemma \ref{lemma-determine-intersections} we see that
$\alpha_j = c_j(\mathcal{F}) \cap [X]$, $j = 0, \ldots, r - 1$
satisfy
$$
\sum\nolimits_{j = 0}^{r - 1} (-1)^jc_1(\mathcal{O}_{P'}(1))^j
\cap (\pi')^*\alpha_j = 0
$$
Pushing this to $P$ and using the remark above as well as
Lemma \ref{lemma-pushforward-cap-c1} we get
$$
\sum\nolimits_{j = 0}^{r - 1}
(-1)^j c_1(\mathcal{O}_P(1))^{j + 1}
\cap \pi^*\alpha_j = 0
$$
By the uniqueness of Lemma \ref{lemma-determine-intersections}
we conclude that
$c_r(\mathcal{E}) \cap [X] = 0$ and
$c_j(\mathcal{E}) \cap [X] = \alpha_j = c_j(\mathcal{F}) \cap [X]$
for $j = 0, \ldots, r - 1$. Hence the lemma holds.
\end{proof}

\end{document}
